\subsection{IV)}

{\noindent\large\begin{align*}\displaystyle F(x, y, z) = \sum(3, 4, 5, 6, 7)\end{align*}}

\subsubsection{Obtener la mínima expresión booleana}
\begin{center}
\begin{tabular}{c|ccc|c}
\rowcolor[HTML]{FFCCC9} 
\cellcolor[HTML]{9AFF99}$m$ & $x$ & $y$ & $z$ & \cellcolor[HTML]{FFFFC7}$F(x,y,z)$ \\ \hline
0                           & 0   & 0   & 0   &                                    \\
1                           & 0   & 0   & 1   &                                    \\
2                           & 0   & 1   & 0   &                                    \\
3                           & 0   & 1   & 1   & 1                                  \\
4                           & 1   & 0   & 0   & 1                                  \\
5                           & 1   & 0   & 1   & 1                                  \\
6                           & 1   & 1   & 0   & 1                                  \\
7                           & 1   & 1   & 1   & 1                                 
\end{tabular}
\end{center}

\begin{center}
    \begin{karnaugh-map}[2][4][1][$z$][$x\qquad{}y$]
        \minterms{3,4,5,6,7}
        \implicant{6}{5}\implicant{3}{7}
    \end{karnaugh-map}
\end{center}

\begin{itemize}[nosep]
    \item El primer bloque ($m_3$ y $m_7$), cambia $x$ por lo que sale y permanecen $y$ y $z$ sin negar.
    \item El segundo bloque ($m_3$, $m_7$, $m_4$ y $m_5$) permanece constante $x$ sin negar y el resto salen.
    \item La función simplificada queda:
\end{itemize}
{\large\begin{align*}F(x,y,z) &= x + z\cdot y\end{align*}}
\subsubsection{Simulación}
\imagen[]{10cm}{./imagenes/ejercicio4Falstad.png}
